\documentclass[11pt,a4paper]{article}

% PAQUETES
\usepackage[T1]{fontenc}%
\usepackage[utf8]{inputenc}%

\usepackage[english]{babel}
\usepackage{amsmath}
\usepackage{amsthm}
\usepackage{amsfonts}
\usepackage[%left=1.54cm,right=1.54cm,top=1.54cm,bottom=1.54cm
    margin=1in,
]{geometry}
\usepackage{xfrac}  
\usepackage{tikz-cd}
\usepackage{enumerate}
\usepackage{amsfonts}
\usepackage{amssymb}
\usepackage{tcolorbox}
\usepackage{rotating}
\usepackage{mathpazo}
% \usepackage{charter}
\usetikzlibrary{babel}
\usepackage{listings}
\usepackage{amssymb}
\usepackage{extarrows}
\usepackage{makeidx}
\usepackage{graphicx}
\usepackage{multirow}
\usepackage{tikz-cd}
\usepackage{tasks}
\usepackage{xcolor}

% Christian
\usepackage{enumitem,etoolbox,titlesec}


%OPERADORES

\DeclareMathOperator{\dom}{dom}
\DeclareMathOperator{\cod}{cod}
\DeclareMathOperator{\id}{id}

\newcommand{\red}[1]{\textcolor{red}{#1}}
\newcommand{\C}{\mathbb{C}}
\newcommand{\im}{\text{Im}}
\newcommand{\R}{\mathbb{R}}
\newcommand{\N}{\mathbb{N}}
\newcommand{\Z}{\mathbb{Z}}
\newcommand{\D}{\mathbb{D}}
\newcommand{\B}{\mbox{Ob}}
\newcommand{\M}{\mbox{Mo}}
\newcommand{\del}{\Delta}
\newcommand{\odel}[1]{\left[#1\right]}
\newcommand{\Hom}[1]{\text{Hom}(#1)}
\newcommand{\adel}[1]{\left\lbrace #1 \right\rbrace}
%\renewcommand{\theequation}{\thesection.\arabic{equation}}
\newcommand{\funcion}[5]{%
{\setlength{\arraycolsep}{2pt}
\begin{array}{r@{}ccl}
#1\colon \hspace{0pt}& #2 & \longrightarrow & #3\\
& #4 & \longmapsto & #5
\end{array}}}

\newcommand{\func}[3]{#1\colon  #2  \to  #3}

\newcommand\restr[2]{{% we make the whole thing an ordinary symbol
  \left.\kern-\nulldelimiterspace % automatically resize the bar with \right
  #1 % the function
  \vphantom{\big|} % pretend it's a little taller at normal size
  \right|_{#2} % this is the delimiter
  }}

%ENTORNOS

% \theoremstyle{theorem}
\newtheorem{teo}{Theorem}[section]
\newtheorem{prop}[teo]{Proposition}
\newtheorem{lem}[teo]{Lemma}
\newtheorem{cor}[teo]{Corollary}

\theoremstyle{definition}
\newtheorem{defi}[teo]{Definition}
\newtheorem{rem}{Remark}[teo]
\newtheorem{exa}{Example}
\newtheorem{eje}{Exercise}
\newtheorem{que}{Question}

\newenvironment{sol}
  {\begin{proof}[\textit{Solution}]}
  {\end{proof}}

\titlelabel{\thetitle.\quad}

\def\contador{}
\graphicspath{{./figures/}}
\newcommand{\qand}{\quad\text{and}\quad}
\usepackage{microtype,parskip}
\setlength{\parindent}{15pt}
\linespread{1.15}
\usepackage{hyperref}
\hypersetup{
    colorlinks=true,  
    allcolors=blue,
    pdfproducer={Christian Chávez},
}

\makeatletter
\@ifclassloaded{exam}{
    \footer{}{\thepage}{}
    \renewcommand{\thequestion}{\bfseries\arabic{question}}
    \renewcommand{\solutiontitle}{\noindent\textit{Solution.}\enspace}
    \unframedsolutions
}{}
\makeatother


\newenvironment{theproof}
{
    \renewcommand{\solutiontitle}{}
    \begin{solution}
    \vspace*{-\baselineskip}
    \begin{proof}
}
{
    \end{proof}
    \end{solution}
    \renewcommand{\solutiontitle}{\noindent\textit{Solution.} }
}

\begin{document}

\def\contador{Lesson 10}
\noindent
\begin{minipage}[c]{0.33\textwidth}
    \includegraphics[scale=0.37]{sello_yachay.png}
\end{minipage}
\begin{minipage}[c]{0.37\textwidth}
    % \centering
    \textbf{\large School of Mathematical and\\ Computational Sciences}\par
    Abstract Algebra
\end{minipage}
~ 
\begin{minipage}[c]{3mm}
    \raggedleft
    \rule[1.5mm]{0.3mm}{15mm}
\end{minipage}
~ 
\begin{minipage}[c]{0.24\textwidth}
    \raggedleft
    Prof. Pablo Rosero\\
    \& Christian Chávez\\
    \contador{}
\end{minipage}

\vspace{1mm}
\noindent\hrulefill

\vspace{3mm}

\section{Ring Homomorphisms and the Fundamental Theorems of Isomorphism}
% idea: define ideales y cocientes
% en otra lección homomorfimos y los teoremas fundamentales de iso.
% Homomorphisms

\subsection{Ring Homomorphisms}

To understand better rings and their properties we have to look at maps between ring that preserve the structure, namely that preserve both addition and multiplication.
This maps are called ring-homomorphisms.

\begin{defi}
    Let \(A\) and \(B\) be rings.
    A map \(\psi\colon A\to B\)  is a \textbf{ring-homomorphism} if, for all \(a,b\in A\),
\begin{enumerate}[label=(\roman*)]
    \item \(\varphi(a+b)=\varphi(a)+\varphi(b)\), and
    \item \(\psi(a\cdot b) = \psi(a)\psi(b)\).
\end{enumerate}
A   bijective ring-homomorphism is called a \textbf{ring-isomorphism}. 
If there is an isomorphism between \(A\) and \(B\), we say these rings are isomorphic, which is denoted by \(A\cong B\).
\end{defi}

Note that condition (i) means \(\psi\) is a group-homomorphism from \((A,+_A)\) to \((B,+_B)\).
This implies that \(\psi(0_A) = 0_B\) since \(\psi(0_A) = \psi(0_A+0_A) = \psi(0_A) + \psi(0_A)\) whence \(\psi(0_A) = 0_B\) by substraction.
A ring-homomorphism from \(A\) to \(A\) is called an \textbf{endomorphism} of \(A\).

\begin{rem}
    If both \(A\) and \(B\) are rings with unity, a ring-homomorphism  \(\psi\colon A\to B\) must in addition satisfy \(\psi(1_A) = 1_B\).
\end{rem}

Given a homomorphism of rings \(\psi\colon A\to B\), the kernel of \(\psi\) is 
\[\operatorname{Ker}\psi =\{x \in A \mid \varphi(x)=0\}.\]
The image of \(\psi\) is \(\im \psi = \psi(A) \).

\begin{eje}
    Prove \(\ker\psi\) is an ideal of \(A\) and \(\im\psi \) is a subring of \(B\). Why is it not necessarily true that \(\im\psi \) is an ideal  of \(B\)?
\end{eje}

\begin{exa}
\begin{enumerate}[label=(\roman*)]
\item If \(B\) is a subring of a ring \(A\), the canonical injection \(\iota \colon B\to A\colon b\mapsto b\) is a ring-homomorphism.
In particular, the identity \(\text{Id}_A\colon A\to A\) is an isomorphism.

\item If \(A\) is a ring and \(I\trianglelefteq A\) an ideal, the canonical projection \(\pi \colon A\to A/I\) is a surjective homomorphism of rings.
Indeed,  we have 
\begin{align*}
    \pi (a_1) + \pi (a_2) & = (a_1 + I) + (a_2 + I)
                    = (a_1 + a_2) + I
                    = \pi (a_1 + a_2), \quad\text{and}\\
    \pi (a_1) \pi (a_2)   & = (a_1 + I)(a_2 + I)
                     = (a_1 a_2) + I 
                     = \pi (a_1 a_2).
\end{align*}
    

In particular, the projection  \(\Z\to \Z/ n\Z\) is a ring homomorphism, called \textbf{reduction modulo} \(n\).

\item If \((A_\lambda)_{\lambda\in \Lambda}\) is a family of rings, the \(\alpha\)th canonical projection \(p_\alpha \colon \prod_{\lambda\in \Lambda} A_\lambda\to A_\alpha\) is a surjective ring-homomorphism.
\item The rings \( 2\mathbb{Z} \) and \( 3\mathbb{Z} \) are not isomorphic. If there were an isomorphism
\(
\psi: 2\mathbb{Z} \to 3\mathbb{Z}
\)
such that \( \psi(2) = 3k \) for some integer  \( k \neq 0 \), then
\[
\psi(4) = \psi(2 + 2) = \psi(2) + \psi(2) = 6k,
\]
but
\[
\psi(4) = \psi(2 \cdot 2) = \psi(2) \cdot \psi(2) = 9k^2.
\]
Thus
\(
6k = 9k^2,
\)
and since \( k \neq 0 \), then \( 2 = 3k \), which is impossible. 
Therefore \( 2\mathbb{Z} \ncong 3\mathbb{Z} \).

\item Let \( A \) be a commutative ring with unity, then the map
\[
\psi: A[x] \to A\quad: \quad p(x) \mapsto  p(0)
\]
is a ring homomorphism that maps a polynomial to its constant term.
In general, the evaluation map \(\psi_\alpha: A[x]\to A\) defined by 
\[
\psi_\alpha\left(a_0 + a_1x + \cdots + a_nx^n\right) = a_0 + a_1\alpha + \cdots + a_n\alpha^n,
\] 
where \(\alpha\in A\) is fixed,  is a homomorphism of \( A[x] \) onto \( A \). 

\item There is no ring homomorphism \( \mathbb{Z}/n\mathbb{Z} \to \mathbb{Z} \) for any \( n \in \Z^+ \).
\end{enumerate}
\end{exa}

\subsection{Properties of Ring Homomorphisms}

\begin{prop}
Let \(\psi\colon A\to B\) be a morphism of rings.
Then
\begin{enumerate}[label=(\roman*)]
\item \( f(-a) = -f(a) \) for all \(a\in A\)
\item \( f(na) = nf(a) \)  for all \(a\in A\) and \(n\in \Z\)
\item \( f(a^n) = f(a)^n \)  for all \(a\in A\) and \(n\in \Z^+\)

\item \(\ker \psi \) is an ideal of \(A\)
\item \(\im \psi\) is a subring of \(B\)
\item \(\psi\) is injective if and only of \(\ker \psi = 0\)
\item if \(a\in A\) is invertible, so is \(\psi(a)\)
\item if \(\mathfrak{p}\) is an ideal of \(B\), then \(\psi^{-1}(\mathfrak{p})\) is an ideal of \(A\).
\item if \(\mathfrak{p}\) is a maximal (prime) ideal of \(B\), then \(\psi^{-1}(\mathfrak{p})\) is a maximal (prime) ideal of \(A\).
\item the composition of ring homomorphisms is a ring homomorphism
\end{enumerate}
\end{prop}

\begin{proof}
    Straight from the definitions and left as an exercise to the student.
\end{proof}

\pagebreak
\subsection{Fundamental Theorems of Isomorphism}

caracterizacion primos y maximales con cocientes




\end{document}